% 学术学位博士研究生中期考核示例。
% 复制本文件后，先修改下列基本信息，再填写各报告部分。
\documentclass[da]{xdupgtm}
% 学号
\newcommand\studentNo{1101110071}
% 论文题目，可使用两行
\newcommand\thesisTitleOne{博士学位论文题目}
\newcommand\thesisTitleTwo{}
% 姓名、一级学科、二级学科、指导教师和学院
\newcommand\authorName{张三}
\newcommand\major{电子科学与技术}
\newcommand\submajor{电磁场与微波技术}
\newcommand\advisorName{李四}
\newcommand\school{电子工程学院}
% 报告日期和选题来源
\newcommand\submitDate{2026年10月8日}
\newcommand\subjectSource{国家自然科学基金}

% 需要电子签名时取消注释，并将图片放入 figures/sign/。
% \renewcommand\zdjsqm{\sign{lisi}}
% \renewcommand\zdjsqmrq{2026年10月8日}
% \renewcommand\zzqm{\sign{lisi}}
% \renewcommand\cyqm{\sign{lisi}\emptysign\sign{lisi}}
% \renewcommand\cyqmrq{2026年10月8日}

\begin{document}
% 选题来源只需在 xzlyj 环境外修改 subjectSource。
\begin{xzlyj}\end{xzlyj}
\begin{yjmbnr}
% 在这里填写学位论文研究目标及研究内容；应与开题报告保持一致。
\end{yjmbnr}
\begin{ywcgz}
% 在这里填写目前已完成的理论、实验研究内容及获得的结论。
\end{ywcgz}
\begin{workplan}
% 第三至第五部分合用一个表格，保留 workplan 外层环境。
\begin{bxfsm}
% 在这里填写与开题报告不相符的情况；没有变化时填写“无”。
% 填写示例：研究目标与开题报告一致，暂无调整。
\end{bxfsm}
\begin{gzhjh}
% 在这里填写下一步工作计划、拟完成的研究内容和需要解决的关键技术。
% 填写示例：完成补充实验与结果分析，完善论文并提交导师审阅。
\end{gzhjh}
\begin{kycg}
% 在这里填写已发表、已录用或拟发表的论文及其他研究成果。
% 填写示例：列出论文题目、作者、期刊或会议、发表状态及与本课题的关系。
\end{kycg}
\end{workplan}
\begin{jsyj}
% 在这里填写指导教师评价意见。
\end{jsyj}
\begin{pyjl}
% 在这里填写中期总结报告评语及“通过”或“不通过”结论。
\end{pyjl}
\end{document}
